% LaTeX snippet for the refinement method implemented in f0_extraction/refine_f0.py.
% Requires: \usepackage{amsmath,amssymb}

\subsection{Elelet-Based Harmonic Refinement}

The manually corrected contour is stored as the first harmonic, denoted
\(H1_i\), at frame \(i\). Refinement is performed in fundamental-frequency
space by converting the first harmonic to its implied \(F0\):
\begin{equation}
    f_i^{(0)} = \frac{H1_i}{2}.
\end{equation}

For each Elelet kernel size \(k \in \mathcal{K}\), the audio signal \(x\) is
transformed into complex Elelet coefficients \(E^{(k)}(x)\). The power
representation used for scoring is
\begin{equation}
    P^{(k)}_{c,m}
    =
    \left|E^{(k)}(x)_{c,m}\right|^{p},
\end{equation}
where \(c\) indexes frequency channels, \(m\) indexes time frames, and \(p\) is
the magnitude exponent.

Let \(A(\cdot)\) denote the auditory-scale frequency transform used by the
Elelet representation, and let \(A^{-1}(\cdot)\) be its inverse. For each input
frame \(i\), the current \(F0\) estimate is mapped to auditory-scale space:
\begin{equation}
    a_i = A\!\left(f_i^{(0)}\right).
\end{equation}
The nearest Elelet frequency bin is
\begin{equation}
    b_i =
    \operatorname*{argmin}_{c}
    \left|a_c - a_i\right|.
\end{equation}

The candidate set is a local neighborhood around this bin:
\begin{equation}
    \mathcal{C}_i =
    \left\{
        b_i + d
        \;:\;
        d \in \{-r_i,\ldots,r_i\}
    \right\},
\end{equation}
where \(r_i\) is the local search radius in auditory bins. Near signal edges,
the implementation reduces this radius to one bin.

Each candidate bin \(c \in \mathcal{C}_i\) corresponds to a candidate
fundamental frequency
\begin{equation}
    \hat{f}_{i,c} = A^{-1}(a_c).
\end{equation}
Candidates below the minimum allowed \(F0\) are discarded.

For a candidate \(\hat{f}_{i,c}\), the maximum number of harmonics considered is
bounded by Nyquist:
\begin{equation}
    H_c =
    \min\left(
        H_{\max},
        \left\lfloor
            \frac{f_s/2}{\hat{f}_{i,c}}
        \right\rfloor
    \right),
\end{equation}
where \(f_s\) is the sampling rate and \(H_{\max}\) is the configured maximum
number of harmonics.

For harmonic \(h\), the target frequency is
\begin{equation}
    f_{h,c} = h \hat{f}_{i,c},
\end{equation}
and its nearest auditory-scale bin in representation \(k\) is
\begin{equation}
    q_{h,c}^{(k)}
    =
    \operatorname*{argmin}_{q}
    \left|
        a_q^{(k)}
        -
        A\!\left(f_{h,c}\right)
    \right|.
\end{equation}

The implementation averages power in a small bin neighborhood around each
harmonic:
\begin{equation}
    \mathcal{B}_{h,c}^{(k)}
    =
    \left\{
        q
        \;:\;
        q_{h,c}^{(k)} - \rho
        \le q \le
        q_{h,c}^{(k)} + \rho
    \right\},
\end{equation}
where \(\rho\) is the harmonic bin radius.

For frame \(i\), representation \(k\), and candidate \(c\), the local harmonic
energy is normalized by the mean frame power:
\begin{equation}
    \bar{P}_{h,c,i}^{(k)}
    =
    \frac{1}{\left|\mathcal{B}_{h,c}^{(k)}\right|}
    \sum_{q \in \mathcal{B}_{h,c}^{(k)}}
    P_{q,m_i^{(k)}}^{(k)},
\end{equation}
\begin{equation}
    Z_i^{(k)}
    =
    \frac{1}{C_k}
    \sum_{q=1}^{C_k}
    P_{q,m_i^{(k)}}^{(k)}
    + \varepsilon.
\end{equation}
Here \(m_i^{(k)}\) is the Elelet frame corresponding to contour frame \(i\),
\(C_k\) is the number of channels for representation \(k\), and
\(\varepsilon\) is a small constant for numerical stability.

The harmonic salience score is then
\begin{equation}
    S_i(c)
    =
    \sum_{k \in \mathcal{K}}
    \sum_{h=1}^{H_c}
    \frac{1}{h^\gamma}
    \frac{
        \bar{P}_{h,c,i}^{(k)}
    }{
        Z_i^{(k)}
    },
\end{equation}
where \(\gamma\) controls the harmonic weighting.

The best candidate is selected by maximum salience:
\begin{equation}
    c_i^\star
    =
    \operatorname*{argmax}_{c \in \mathcal{C}_i}
    S_i(c).
\end{equation}

When enabled, a quadratic sub-bin correction is applied to the winning
candidate using its neighboring scores:
\begin{equation}
    \Delta_i
    =
    \frac{1}{2}
    \frac{
        S_i(c_i^\star - 1) - S_i(c_i^\star + 1)
    }{
        S_i(c_i^\star - 1)
        -
        2S_i(c_i^\star)
        +
        S_i(c_i^\star + 1)
    },
\end{equation}
\begin{equation}
    \tilde{c}_i^\star
    =
    \operatorname{clip}
    \left(
        c_i^\star + \Delta_i,
        1,
        C - 2
    \right).
\end{equation}
If sub-bin interpolation is disabled, then \(\tilde{c}_i^\star = c_i^\star\).

The refined fundamental frequency is obtained by interpolating the auditory
frequency axis at the selected bin and mapping back to Hz:
\begin{equation}
    \hat{f}_i
    =
    A^{-1}
    \left(
        \operatorname{interp}
        \left(
            \tilde{c}_i^\star,
            \{a_c\}_{c=1}^{C}
        \right)
    \right).
\end{equation}

If median smoothing is enabled, it is applied only inside contiguous voiced
regions:
\begin{equation}
    \hat{f}_i^{\,\mathrm{smooth}}
    =
    \operatorname{median}
    \left(
        \hat{f}_{i-w},
        \ldots,
        \hat{f}_{i},
        \ldots,
        \hat{f}_{i+w}
    \right),
\end{equation}
where \(2w+1\) is the median filter length.

Finally, the refined contour is converted back to first-harmonic space for
output:
\begin{equation}
    H1_i^{\mathrm{refined}}
    =
    2\hat{f}_i^{\,\mathrm{smooth}}.
\end{equation}

